% Lösungen Einheit 1 von 4 – Terme aufstellen und berechnen (zusammenbau v0.9)
\einheitenkopf{Einheit 1 von 4 · Terme aufstellen und berechnen}

\erg{6}{a) $4x - 3$ \quad b) $4x - 5$ \quad c) $4x - 6$ \quad d) $4x - 7$ \quad e) $3x + 8$}
\erg{7}{a) $5x$ (auch $5 \cdot x$) \quad b) $4a + 3b$ \quad c) $x + 9 + x + 9 = 2x + 18$ \quad d) $2 \cdot (3x + 2)$}
\erg{8}{a) Dreifaches $3x$, Differenz $3x - 7$, verdoppelt mit Klammer: $2 \cdot (3x - 7)$ \quad b) Sechsfaches $6x$, dann um 2 vermindert: $6x - 2$}
\erg{9}{a) $4 \cdot 5 - 3 = 17$ \quad b) $2 \cdot (-3) + 9 = 3$ \quad c) $10 - 3 \cdot (-2) = 16$}
\erg{10}{a) Zum Beispiel: Ein Kinobesuch kostet 15 € Eintritt, dazu kommen $x$ Getränke zu je 3 €; $x$ ist die Zahl der Getränke. Einsetzen: $15 + 3 \cdot 2 = 21$, der Besuch kostet 21 €. \quad b) Zum Beispiel: Lara hat 20 € und kauft $x$ Hefte zu je 2 €; $x$ ist die Zahl der Hefte. Für $x = 4$: $20 - 2 \cdot 4 = 12$, ihr bleiben 12 €. \quad c) Zum Beispiel: Ein Fahrradverleih nimmt 10 € Grundgebühr und 4 € je Stunde; $x$ ist die Zahl der Stunden. Für $x = 3$: $10 + 4 \cdot 3 = 22$, drei Stunden kosten 22 €.}
\erg{11}{a) Richtig: $x - 9$. Mia hat die Reihenfolge vertauscht; „vermindert um 9“ heißt, die 9 wird von $x$ abgezogen. \quad b) Richtig. Verdoppelt wird die ganze Summe, deshalb braucht sie die Klammer: $2 \cdot (x + 7)$. \quad c) Zeile 2 ist falsch: Verdoppelt wird die ganze Summe, die Klammer fehlt. Richtig: $2 \cdot (x + 6)$.}
\erg{12}{a) Die Klammer sorgt dafür, dass die ganze Summe verdoppelt wird, nicht nur $x$; Probe mit $x = 3$: $2 \cdot (3 + 5) = 16$, aber $2 \cdot 3 + 5 = 11$ – die Terme haben verschiedene Werte. \quad b) (1) falsch, z. B. $x = 10$: $10 - 3 = 7$, aber $3 - 10 = -7$. (2) wahr, denn ohne Klammer würde nur das erste Glied verdoppelt. (3) falsch, denn $2 \cdot (x + 4) = 2x + 8$ ist für jede Zahl $x$ um 4 größer als $2x + 4$. \quad c) Nein; die 2 muss die ganze Summe treffen, dafür braucht die Summe eine Klammer: $2 \cdot (x + 6) = 2x + 12$, und das ist nicht $2x + 6$.}
\erg{13}{a) $4 \cdot (x + 6)$ \quad b) Oben $-$, links der Ast $2 \cdot x$, rechts 7. \quad c) Oben $\cdot$, links der Ast $x - 2$, rechts 5.}

13b)

\termbaum{-}{\tb{\cdot}{2}{x}}{7}


13c)

\termbaum{\cdot}{\tb{-}{x}{2}}{5}

\erg{14}{a) Term aufstellen: $11 + 0{,}30n$; einsetzen: $11 + 0{,}30 \cdot 150 = 56$ €. Ein Monat mit 150 Kilowattstunden kostet 56 €. \quad b) Nein; Term aufstellen: $4{,}50 + 2{,}50d$; einsetzen: $4{,}50 + 2{,}50 \cdot 9 = 27$ €, das sind 2 € mehr als 25 €. \quad c) $30 + 45t$; für 3 Stunden: $30 + 45 \cdot 3 = 165$ €}
